%% sections/abstract.tex
\begin{center}
\begin{minipage}{0.88\textwidth}
\small
\noindent\textbf{Abstract.}
We replicate \citet{dufwenberg2005bubbles} computationally and extend
the design with the expected-utility agent of
\citet{lopezlira2025llm}.  Six DLM traders with an endogenous
experience counter reproduce the original finding: the bubble dies
by round~4.  Replacing DLM traders with heterogeneous utility agents
that carry risk preferences, valuation biases, and a pairwise trust
channel yields a qualitatively similar attenuation pattern under an
algorithmic belief blend (Plan~I).  We then substitute the blend with
a language-model call at every period boundary, comparing a prompt
that includes the closed-form utility expression (Plan~II) against one
that names only the risk label (Plan~III).  Plans~II and~III produce
indistinguishable bubble statistics, indicating that the model's
internal representation of risk types is robust to whether the
functional form is stated explicitly.  All three plans share the same
market substrate and seeded PRNG, so differences are attributable
to the belief channel alone.

\medskip
\noindent\textbf{JEL:} C63, C92, D84, G12.
\quad\textbf{Keywords:} asset bubbles, experience, expected utility,
language models, belief updating.
\end{minipage}
\end{center}
\bigskip
